% Partie : sets-functions-relations
% Chapitre : sets
% Section : pairs-and-products

\documentclass[../../../include/open-logic-section]{subfiles}

\begin{document}

\olfileid[fr]{sfr}{set}{pai}
\olsection[Couples, n-uplets et produits cartésiens]{Couples, $n$-uplets et produits cartésiens}

\begin{explain}
L'extensionnalité implique que les éléments d'un ensemble ne
sont pas ordonnés. Pour représenter un ordre, nous utilisons
des \emph{couples}, ou \emph{paires ordonnées}, $\tuple{x, y}$.
Dans une paire non ordonnée $\{x, y\}$, l'ordre n'importe pas :
$\{x, y\} = \{y, x\}$. Il importe dans un couple :
si $x \neq y$, alors $\tuple{x, y} \neq \tuple{y, x}$.

Comment représenter les couples en théorie des ensembles ?
Il est essentiel que deux couples soient égaux si et seulement
s'ils ont la même première composante et la même seconde
composante, c'est-à-dire :
\[
  \tuple{a, b}= \tuple{c, d}\text{ si et seulement si }a = c \text{ et }b=d.
\]
La définition de Wiener--Kuratowski permet de définir ainsi
les couples en théorie des ensembles.
\end{explain}

\begin{defn}[Couple]\ollabel{wienerkuratowski}
	$\tuple{a, b} = \{\{a\}, \{a, b\}\}$.
\end{defn}

\begin{prob}
	À l'aide de la \olref[sfr][set][pai]{wienerkuratowski},
	démontrer que $\tuple{a, b}= \tuple{c, d}$ si et seulement si
	$a = c$ et $b=d$.
\end{prob}

\begin{explain}
Une fois les couples définis, nous pouvons les utiliser pour
définir d'autres ensembles. Nous avons parfois besoin de suites
ordonnées de plus de deux objets : des \emph{triplets}
$\tuple{x, y, z}$, des \emph{quadruplets} $\tuple{x, y, z, u}$,
et ainsi de suite. Un triplet peut être représenté par un couple
dont la première composante est elle-même un couple :
$\tuple{x, y, z}$ est $\tuple{\tuple{x, y},z}$.
De même, le quadruplet $\tuple{x,y,z,u}$ est
$\tuple{\tuple{\tuple{x,y},z},u}$, et ainsi de suite.
On parle en général de \emph{$n$-uplets ordonnés}
$\tuple{x_1, \dots, x_n}$.

Certains ensembles de couples, ou plus généralement de
$n$-uplets ordonnés, nous seront utiles.
\end{explain}

\begin{defn}[Produit cartésien]
Étant donnés deux ensembles $A$ et $B$, leur \emph{produit cartésien}
$A \times B$ est défini par
\[
  A \times B = \Setabs{\tuple{x, y}}{x \in A \text{ et } y \in B}.
\]
\end{defn}

\begin{ex}
Si $A = \{0, 1\}$ et $B = \{1, a, b\}$, leur produit est
\[
A \times B = \{ \tuple{0, 1}, \tuple{0, a}, \tuple{0, b},
    \tuple{1, 1}, \tuple{1, a}, \tuple{1, b} \}.
\]
\end{ex}

\begin{ex}
Si $A$ est un ensemble, le produit de $A$ avec lui-même, $A \times A$,
se note aussi~$A^2$. C'est l'ensemble de \emph{tous} les couples
$\tuple{x, y}$ tels que $x, y \in A$.
L'ensemble de tous les triplets $\tuple{x, y, z}$ est $A^3$,
et ainsi de suite. On peut en donner une définition par récurrence :
\begin{align*}
  A^1 & = A\\
  A^{k+1} & = A^k \times A
\end{align*}
\end{ex}

\begin{prob}
Énumérer tous les !!{element}s de $\{1, 2, 3\}^3$.
\end{prob}

\begin{prop}\ollabel{cardnmprod}
Si $A$ a $n$ !!{element}s et $B$ a $m$ !!{element}s, alors
$A \times B$ a $n\cdot m$ éléments.
\end{prop}

\begin{proof}
Pour chaque !!{element}~$x$ de~$A$, il y a $m$ !!{element}s
de la forme $\tuple{x, y} \in A \times B$.
Posons $B_x = \Setabs{\tuple{x, y}}{y \in B}$.
Puisque $x_1 \neq x_2$ entraîne
$\tuple{x_1, y} \neq \tuple{x_2, y}$, on a
$B_{x_1} \cap B_{x_2} = \emptyset$.
Si $A = \{x_1, \dots, x_n\}$, alors
$A \times B = B_{x_1} \cup \dots \cup B_{x_n}$,
qui possède donc $n\cdot m$ !!{element}s.

Pour le visualiser, disposons les !!{element}s de~$A \times B$
dans un tableau :
\[
\begin{array}{rcccc}
  B_{x_1} = & \{\tuple{x_1, y_1} & \tuple{x_1, y_2} & \dots & \tuple{x_1, y_m}\}\\
  B_{x_2} = & \{\tuple{x_2, y_1} & \tuple{x_2, y_2} & \dots & \tuple{x_2, y_m}\}\\
  \vdots & & \vdots\\
  B_{x_n} = & \{\tuple{x_n, y_1} & \tuple{x_n, y_2} & \dots & \tuple{x_n, y_m}\}
\end{array}
\]
Les $x_i$ étant tous distincts, de même que les $y_j$,
les couples de ce tableau sont deux à deux distincts,
et il y en a $n\cdot m$.
\end{proof}

\begin{prob}
Montrer par récurrence sur~$k$ que, pour tout $k \ge 1$,
si $A$ a $n$ !!{element}s, alors $A^k$ a $n^k$ !!{element}s.
\end{prob}

\begin{ex}
Si $A$ est un ensemble, un \emph{mot} sur~$A$ est une suite
finie d'!!{element}s de~$A$. Une telle suite peut être considérée
comme un $n$-uplet d'!!{element}s de~$A$. Par exemple, si
$A = \{a, b, c\}$, la suite «~$bac$~» peut être considérée
comme le triplet~$\tuple{b, a, c}$.
Les mots, c'est-à-dire les suites de symboles, jouent un rôle
essentiel en informatique. Par convention, on considère les
!!{element}s de~$A$ comme des suites de longueur~$1$,
et $\emptyset$ comme la suite de longueur~$0$.
Avec ce codage, l'ensemble des représentations de \emph{tous}
les mots sur~$A$ s'écrit alors\footnote{Précision de l'édition
française : pour un alphabet arbitraire, ce codage peut identifier
des mots de longueurs différentes. Par exemple, si le symbole
$\emptyset$ appartient à $A$, le mot vide et le mot d'une seule
lettre $\emptyset$ ont la même représentation. Pour représenter
les mots sans ambiguïté, on associe en général à chaque
représentation la longueur du mot. La formule de l'original
est conservée ci-dessous comme description des représentations
non munies de leur longueur.}
\[
A^* = \{\emptyset\} \cup A \cup A^2 \cup A^3 \cup \dots
\]
\end{ex}

\end{document}
